\documentclass[11pt,a4paper]{article}
\usepackage[T1]{fontenc}
\usepackage[utf8]{inputenc}
\usepackage{lmodern,amsmath,amssymb,amsthm,booktabs,microtype}
\usepackage[margin=29mm,top=25mm,bottom=26mm]{geometry}
\usepackage[colorlinks=true,linkcolor=black,citecolor=black,urlcolor=blue!50!black]{hyperref}
\usepackage{xcolor}
\hypersetup{pdftitle={A Lean 4 formalization of the Bhargava--Hanke 290 theorem},pdfauthor={Chang Su}}
\newtheorem{theorem}{Theorem}
\newtheorem{lemma}[theorem]{Lemma}
\newcommand{\ZZ}{\mathbb Z}
\newcommand{\NN}{\mathbb N}
\newcommand{\rep}{\operatorname{Rep}}
\newcommand{\code}[1]{\texttt{#1}}
\title{A Lean 4 formalization of the\\ Bhargava--Hanke 290 theorem}
\author{Chang Su}
\date{27 November 2026}
\begin{document}
\maketitle
\begin{abstract}
The Bhargava--Hanke 290 theorem characterizes universal positive-definite
quadratic forms with integer coefficients by representation of 29 specified
integers. We give a Lean 4 proof for forms of arbitrary finite rank over
$\ZZ$. Representations of the test integers force the ambient form to
contain one of 34 ternary models as a pullback. The remaining argument
combines escalation, explicit maps between lattices, and representation
theorems for squarefree integers. For a typical quaternary model a finite
table chooses a shift for each admissible residue class. A ternary
representation theorem turns that choice into a representation of every
squarefree integer in the class above an explicit threshold. We describe
two such branches and the use of the ambient hypotheses for exceptions.
The accompanying source and proof-term archive support a replay of the two
public theorems in an empty Lean environment.
\end{abstract}

\section{The theorem and its formal statement}
A quadratic form $Q$ on $\ZZ^d$ represents $n$ if $Q(x)=n$ for some
$x\in\ZZ^d$. Write $\rep_Q(n)$ for this assertion. The coefficients of
the polynomial $\sum_{i\leq j}a_{ij}x_ix_j$ are integers. In particular,
an odd coefficient of $x_ix_j$ is allowed, even though the corresponding
entry of its symmetric Gram matrix is then a half-integer.
Let
\[
\begin{split}
S=\{&1,2,3,5,6,7,10,13,14,15,17,19,21,22,23,26,\\
    &29,30,31,34,35,37,42,58,93,110,145,203,290\}.
\end{split}
\]
\begin{theorem}[Bhargava--Hanke {\cite[Theorem~1]{BH}}]\label{thm:main}
Let $d\in\NN$ and let $Q$ be a positive-definite quadratic form on
$\ZZ^d$ with integer coefficients. If $\rep_Q(s)$ for all $s\in S$, then
$\rep_Q(n)$ for every positive integer $n$.
\end{theorem}
In Lean the ambient type is \code{QuadraticForm} over the integers on
the module of functions from $\operatorname{Fin}(d)$ to $\ZZ$. The
hypothesis \code{Q.PosDef} means that $Q(x)>0$ whenever $x\ne 0$.
The predicates \code{Represents}, \code{Checks}, and \code{Universal}
in \code{Freely290.Core} respectively express representation, representation
on a finite set, and representation of all positive integers. The public
declaration \code{Freely290.strong290} proves Theorem~\ref{thm:main}.
Since every member of $S$ lies in $[1,290]$, testing this entire interval
suffices as well. This corollary is \code{Freely290.weak290}. The reverse
implication of each test follows immediately from universality.

The original preprint also proves that every member of $S$ is indispensable
and classifies the universal quaternary forms \cite[Theorems~2 and~4]{BH}.
The two Lean declarations here concern the universality criterion. A related
LeanEval problem, submitted by Bolton Bailey/Project Numina, gives a
matrix-based statement of the same mathematical problem \cite{LeanEval}.
On 14 November 2026, LeanEval recorded a successful solution by the
submitter \code{plby} \cite{plby}. The public result records a comparator
pass and identifies the submission as private at the time. The checked declarations in this article use mathlib's integral
quadratic-form interface \cite{Mathlib}.

\section{Forcing a finite list of subforms}
Suppose that $Q$ represents all of $S$. A convenient way to keep track of
chosen vectors is an integral linear map $f\colon\ZZ^k\to\ZZ^d$ with
$A=Q\circ f$. Every representation $A(y)=n$ then supplies the witness
$f(y)$ for $Q$. This implication holds for maps of arbitrary rank, which
matters when adjoining a witness leaves the rational span unchanged.

To see the first escalation explicitly, choose $u,v$ of norms $1,2$.
Their pullback is $x^2+cxy+2y^2$. Replacing $v$ by $-v$ makes $c\geq0$.
Nonnegativity of the pullback gives $c^2\leq8$, so $c\in\{0,1,2\}$.
In the last case $v-u$ has norm $1$ and is orthogonal to $u$.
The resulting binary forms and their least omitted positive integers
(their \emph{truants}) are
\[
\begin{array}{c|ccc}
 A & x^2+y^2 & x^2+xy+2y^2 & x^2+2y^2\\ \hline
 \text{truant} &3&3&5.
\end{array}
\]
Both truants belong to $S$. Thus $Q$ supplies a third vector with the
required norm, even when its binary pullback cannot represent that norm.

At ternary rank the coefficient bounds give 351 parameter choices.
The admissibility tests retain 173. For each retained choice a certificate
gives points for the integers preceding its truant and bounds the search
for a point of the missing norm. Integral coordinate changes reduce the
list to 35 representatives. One of these already contains a retained
model on an index-two sublattice. Pullback along that inclusion gives a
cover by 34 ternary models. Only the forward implication about
representation is needed for this last reduction.

Denote the models by $T_i$. The obligation for model $i$ is an
\emph{ambient} assertion: if $Q$ is nonnegative, $Q$ represents $S$, and
$Q\circ f=T_i$ for some integral $f$, then $Q$ is universal. This
distinction explains how a ternary model that misses its truant can
still be part of the proof. One may adjoin a witness supplied by $Q$ for
a member of $S$, classify the resulting pullback, and pass its
representations back to $Q$. One may also embed a previously treated model
and invoke its ambient assertion. In the assembly for models 26 and 28,
both transfers use conclusions already established for the other 32
models. The final assembly is in \code{FourFreelySeeds} and its successors.

\begin{table}[ht]
\centering
\begin{tabular}{lrrrr}
\toprule
Ternary model index & 23 & 25 & 27 & 32\\
Quaternary representatives in its branch & 963 & 848 & 950 & 481\\
\bottomrule
\end{tabular}
\caption{Four case lists in the final reduction. Indices and counts refer
to the formal development. A representative in a list may be discharged
by transfer from another representative.}\label{tab:cases}
\end{table}
For example, of the 950 targets for model 27, 800 have direct proofs.
The other 150 have checked integral maps from already treated forms.
The finite list establishes coverage of possible pullbacks. Universality
for such a pullback, or for its ambient form, still requires an argument
valid for arbitrarily large integers.

\section{From residues to integer witnesses}
Every positive integer has the form $a b^2$ with $a$ positive and
squarefree. If $Q(x)=a$, then $Q(bx)=a b^2$. It therefore suffices
throughout this section to produce points of squarefree norm.
The following elementary lemma isolates one way that finitely many
checks give infinitely many such points.

\begin{lemma}[Shifts in a represented residue class]\label{lem:shift}
Suppose $A$ and $B$ are integral quadratic forms and an integral linear
map $E$ satisfies
\[
 A(E(y,z))=s B(y)+d z^2,\qquad s>0,\quad d\geq0 .
\]
Let $m>0$, and let $G$ be a set of residues modulo $m$ such that $B$
represents every nonnegative integer in those classes. Fix a multiple
$P>0$ of $sm$ and a threshold $H$. For every residue $r$ modulo $P$
that occurs for a positive squarefree integer, suppose there are integers
$z_r,g_r$ with $g_r\bmod m\in G$ and
\[
 d z_r^2\leq H,\qquad
 r-d z_r^2\equiv s g_r \pmod{sm}.
\]
Then $A$ represents every squarefree $n\geq H$ with $n>0$.
\end{lemma}
\begin{proof}
Choose the row $r=n\bmod P$. The congruence implies that
$t=(n-dz_r^2)/s$ is an integer congruent to $g_r$ modulo $m$.
The bound gives $t\geq0$. Choose $y$ with $B(y)=t$.
Substitution gives $A(E(y,z_r))=n$.
\end{proof}

Thus a table entry specifies how to construct a witness from a ternary
representation theorem. That theorem must supply an actual integral
point. A local congruence on its own supplies no such point.
The implementation in \code{FinitePeriodicTernaryTail} permits the
embedding and its coefficients to vary by row. Other branches use
rational transformations. If integral matrices $M,N,T$ and an integer
$e>0$ satisfy $T^{\mathsf t}MT=e^2N$, a vector $x$ for $N$ gives a
vector $Tx/e$ for $M$ whenever $Tx\equiv0\pmod e$ coordinatewise.
Consequently a transfer certificate checks both the identity and
coverage of the input vector classes on which some transformation is
integral. This use of residue classes to transfer between ternary forms
has a precedent in Oh's method \cite[Lemma~2.2 and Theorem~2.3]{Oh}.

\subsection{A four-variable example that can be checked by hand}
Consider $F=x^2+y^2+z^2+7w^2$. For squarefree $n$, one has $4\nmid n$.
By the three-square theorem, $n$ is a sum of three squares if
$n\not\equiv7\pmod8$. Take $w=0$ then. If $n\equiv7\pmod8$ and $n>28$,
take $w=2$. Now $n-28>0$ is $3$ modulo $8$, so it too is a sum of
three squares. The remaining values $7,15,23$ are represented by taking
$w=1$ and using $0$, $2^2+2^2$, and $4^2$, respectively. This proves
universality by homogeneity. For example the same shift gives
$391=11^2+11^2+11^2+7\cdot2^2$. The formal argument treats a family
of fourth coefficients in \code{ThreeSquaresTail}. Bhargava and Hanke
already use three-square sublattices and squarefree reduction
\cite[\S4.1]{BH}.

\subsection{A periodic tail and an ambient exception}
Target 113 in the branch for model 27 has polynomial
\[
 x^2+2y^2+5z^2+29w^2-xz-7xw-yz-9yw+9zw .
\]
The certificate uses
period 260 and the guard $r\not\equiv0\pmod4$.
A squarefree input satisfies that guard. It leaves 195 possible rows.
The tail threshold is $872$, and the allowed shifts are
\[
                   w\in\{1,2,3,4,5,6,11,12\}.
\]
For a chosen $w$ set $k=37n-224w^2$. The inequality
$224\cdot12^2=32256<37\cdot872$ ensures $k>0$ throughout the tail.
A checked ternary criterion and transfers supply a representation of
$37k$. For instance, when $n\equiv3\pmod4$, take $w=1$.
Then $k\equiv35\pmod{37}$ is a quadratic nonresidue and
$37k\equiv3\pmod4$, as required by one of the two ternary criteria.
The other supplies admissible residues modulo 65. The joint test uses
$\operatorname{lcm}(4,65)=260$. A coset-lifting identity then gives an integral point of the
target of norm $n$. The finite test covers all the guarded rows.
Separate checked representations and ambient exceptions handle the
initial interval $1\le n<872$. The file \code{Candidate27ShortBinaryFrom113}
combines the two ranges into the assertion about $Q$.

Target 847 in the branch for model 25 illustrates what an ambient
exception means. Its polynomial is $x^2+2y^2+5z^2+29w^2-xz-3yz$.
Its finite computation leaves seven values to discharge:
\[
                 145,\ 203,\ 290,\ 4205,\ 5887,\ 8410,\ 10933.
\]
The last four are $29^2$ times $5,7,10,13$, respectively, and cannot be
squarefree. The first three belong to $S$, so the original
hypothesis supplies points in $Q$ for them. The certificate therefore
closes this branch without requiring points of those three norms in
the quaternary pullback. The corresponding declaration is in
\code{Candidate25BinaryWindow847}.

\section{Relation to the original proof}
Both arguments use escalation to extract low-rank forms from an
arbitrary ambient lattice. Both use subforms, finite computations, and
arithmetic information valid for all sufficiently large inputs. In their
treatment of difficult quaternary lattices, Bhargava and Hanke study
the weight-two theta series: its coefficients count representations.
Bounds on the cuspidal part, estimates for the Eisenstein coefficients,
and explicit local-density calculations leave a finite range to check
\cite[\S\S4.2--4.4]{BH}. Their finite check also uses split local
covers and a compact table of ternary representations
\cite[\S4.3.2]{BH}. The formal branches described above obtain the
large-input implication from ternary representations, congruence
coverage, and coordinate maps. They record enough arithmetic data for
Lean to prove that the witnesses lift. This provides a route to the same
universality criterion using a different proof of the required
quaternary implications.

\section{Checking the artifact}
The public entry module is \code{Freely290.Completed290}. The release
pins Lean~4.35.0-rc3 \cite{Lean} and mathlib commit
\code{2a885768dae569d938bb9ff3474da6a8753bb90a}.
The source manifest records the SHA-256 digests of 4,803 local Lean
files in the public dependency closure. The archived proof-term stream
also contains imported declarations needed to type-check its two roots.
It contains 374,653 declarations, including 255,769 theorem declarations.
These counts include dependencies from Lean and mathlib.

The recorded run built the targets, exported their dependencies, and
replayed the serialized declarations with Lean's kernel in an empty
environment. The three axioms in the export are \code{propext},
\code{Classical.choice}, and \code{Quot.sound}. The audit found zero
unsafe or partial declarations. Search programs can suggest shifts,
matrices and vectors. Their suggestions enter the proof through checked
identities, congruences and coverage statements. The saved stream lets a
reader check the resulting terms without running the search.

With Elan, the specified Lean toolchain, and the compressed proof file
in the current directory, the following commands replay the archive on
a POSIX system with \code{/dev/stdin}. The pipefail setting is available
in Bash and Zsh.
\begin{verbatim}
elan toolchain install leanprover/lean4:v4.35.0-rc3
set -o pipefail
gzip -dc freely290-proof.ndjson.gz | \
  elan run leanprover/lean4:v4.35.0-rc3 \
  leanchecker --from-export /dev/stdin
\end{verbatim}
The SHA-256 digest of the compressed stream is\\
{\small\ttfamily 076e02f482d7fbc353bc11b86013e740c72e24245e75b60b62f823daafb641e7}.
The source archive gives instructions for building the Lean files against
the pinned dependencies. The kernel replay checks the exported terms;
the source manifest and the build record identify the source development
that produced them.

\paragraph{Availability and licensing.}
The source, proof terms, and reproducibility instructions accompany
this preprint. The version DOI reserved for the deposit is \href{https://doi.org/10.5281/zenodo.24570786}{10.5281/zenodo.24570786}.
The manuscript is licensed under
\href{https://creativecommons.org/licenses/by/4.0/}{CC BY 4.0}.
The author's project code is licensed under
\href{https://www.apache.org/licenses/LICENSE-2.0}{Apache-2.0}.
Third-party components retain their own licenses.

\paragraph{Use of AI tools.}
OpenAI Codex assisted in developing the Lean proofs, generating
certificates, and preparing the manuscript. The archived verification
record identifies the declarations checked by the Lean kernel.

\begin{thebibliography}{9}
\bibitem{BH} M. Bhargava and J. Hanke.
\emph{Universal quadratic forms and the 290-Theorem}. Preprint.
\href{https://math.stanford.edu/~vakil/files/290-Theorem-preprint.pdf}{Author manuscript}.
\bibitem{LeanEval} B. Bailey/Project Numina.
\emph{The 290 theorem}. LeanEval problem statement, revision~1.
\href{https://github.com/leanprover/lean-eval/blob/f03b744959aac50ab1dd8ee2d4dd4397df5d4e6f/LeanEval/NumberTheory/TwoNinetyTheorem.lean}{Lean source, revision f03b744}.
\bibitem{plby} \code{plby}. LeanEval submission 2186, \emph{The 290 theorem}.
14 November 2026.
\href{https://github.com/leanprover/lean-eval-submissions/issues/2186}{Submission and comparator result}.
\href{https://github.com/leanprover/lean-eval-submissions/commit/48c742b7c39ecb339fe37ec91b7f8ad81fc9e645}{Permanent result record}.
\bibitem{Lean} L. de Moura and S. Ullrich.
\emph{The Lean 4 Theorem Prover and Programming Language}.
CADE 28, 625--635 (2021).
\href{https://doi.org/10.1007/978-3-030-79876-5_37}{doi:10.1007/978-3-030-79876-5\_37}.
\bibitem{Mathlib} The mathlib Community.
\emph{The Lean Mathematical Library}.
CPP 2020, 367--381 (2020).
\href{https://doi.org/10.1145/3372885.3373824}{doi:10.1145/3372885.3373824}.
\bibitem{Oh} B.-K. Oh.
\emph{Regular positive ternary quadratic forms}.
\emph{Acta Arithmetica} \textbf{147}(3), 233--243 (2011).
\href{https://doi.org/10.4064/aa147-3-3}{doi:10.4064/aa147-3-3}.
\end{thebibliography}
\end{document}
